\chapter{Произведение KO2, семейный пфаффиан и предел нормированного веса}

Предыдущая глава нашла числа $1/10$, $1/20$ и $1/14$ внутри допустимого
окна слабого барьера. Теперь необходимо проверить, производит ли их
обычный фермионный функциональный интеграл, а не только нормированный
матричный детерминант.

\section{Произведение внешней и внутренней геометрий}

Пусть внешняя евклидова спиновая геометрия имеет KO-степень четыре, а
конечный Real-модуль проекта --- KO-степень шесть. Для чётного
произведения используются
\begin{align}
 \mathcal H&=\mathcal H_M\widehat\otimes\mathcal H_F,\\
 D&=D_M\otimes1+\gamma_M\otimes D_F,\\
 \gamma&=\gamma_M\otimes\gamma_F,
\end{align}
с подходящим произведением вещественных структур. KO-степени складываются:
\begin{equation}
 4+6=10\equiv2\pmod8.
 \label{eq:v6-product-ko2-sum}
\end{equation}
Именно суммарная степень два допускает физический киральный вещественный
интеграл в форме пфаффиана. Конструкция произведений и сложение
KO-степеней даны у Дабровского--Доссены, а роль внутренней степени шесть
в четырёхмерном произведении --- у Конна; см.
\path{arXiv:1011.4456} и \path{hep-th/0608226}. Классификация конечных
фермионных интегралов приведена у Барретта, \path{arXiv:2403.18428}.

Следовательно, прежнее препятствие внутренней KO6-геометрии снято на
кинематическом уровне: корректный product-KO2 пфаффиан допустим.

\section{Явная семейная антисимметричная форма}

На активной Real-паре возьмём
\begin{equation}
 M_R=I_{15}\otimes R,
 \qquad
 A_R^{\rm act}=
 \begin{pmatrix}
 0&M_R\\
 -M_R^T&0
 \end{pmatrix}.
 \label{eq:v6-ko2-active-family-form}
\end{equation}
Её размерность равна $90$. Добавим независимый от $R$ стандартный
симплектический блок размерности $210$, чтобы получить
$A_R\in M_{300}(\mathbb R)$. Тогда
\begin{equation}
 |\operatorname{Pf}A_R|
 =|\det M_R|
 =\det(R)^{15}.
 \label{eq:v6-ordinary-family-pfaffian}
\end{equation}
При $R\mapsto URU^T$, $U\in O(3)$, модуль выражения сохраняется, поэтому
семейный блок имеет требуемый инвариантный тип.

Таким образом, обычный березинский интеграл даёт
\begin{equation}
 \Gamma_F(R)=-\log|\operatorname{Pf}A_R|
 =-15\log\det R+\text{const}.
 \label{eq:v6-ordinary-pfaffian-barrier}
\end{equation}
Это не коэффициент $1/20$, а коэффициент $15$.

\section{Почему нормировка квадратичной формы не делит показатель}

Можно попытаться вставить нормированный след непосредственно в
фермионное действие:
\begin{equation}
 S_F=\frac1{2\cdot300}\psi^TA_R\psi.
\end{equation}
Однако гауссов интеграл равен
\begin{equation}
 \operatorname{Pf}\!\left(\frac1{300}A_R\right)
 =300^{-150}\operatorname{Pf}(A_R).
\end{equation}
Множитель $300^{-150}$ не зависит от $R$. Поэтому для любых $R,R_0$
\begin{equation}
 \log\frac{|\operatorname{Pf}(A_R/300)|}
 {|\operatorname{Pf}(A_{R_0}/300)|}
 =15\log\frac{\det R}{\det R_0}.
 \label{eq:v6-pfaffian-scale-cancels}
\end{equation}
Нормировка квадратичного действия меняет только постоянную свободной
энергии и не превращает пятнадцать копий в долю одной копии.

Внешние моды также умножают, а не делят показатель. Для $k$ независимых
внешних мод получаем $15k\log\det R$. Поэтому само произведение KO2 не
выбирает знаменатель $300$ или $210$.

\begin{theorem}[Обычный пфаффиан не даёт слабого барьера]
Product-KO2 геометрия разрешает семейный пфаффиан, но стандартный
березинский интеграл на пятнадцатикратном физическом секторе индуцирует
$-15\log\det R$. Этот коэффициент значительно превышает границу
$17/168$ и уничтожает локальную отрицательную моду кристаллизации.
\end{theorem}

\section{Где действительно возникает число 1/20}

Детерминант Фугледе--Кадисона использует нормированный след:
\begin{equation}
 \log\Delta_{\tau_{300}}(A_R)
 =\tau_{300}(\log|A_R|)
 =\frac1{300}\log|\det A_R|.
\end{equation}
Для пфаффиана соответствующая интенсивная величина равна
\begin{equation}
 \frac1{300}\log|\operatorname{Pf}A_R|
 =\frac1{20}\log\det R.
 \label{eq:v6-fk-intensive-one-twentieth}
\end{equation}
Это точная математическая формула, но она является нормированной
свободной энергией на одно внутреннее состояние, а не значением обычного
конечномерного фермионного интеграла. Различие между обычным и
нормированным детерминантами обсуждается в
\path{arXiv:1107.1059}.

Чтобы использовать \eqref{eq:v6-fk-intensive-one-twentieth} физически,
необходимо единым принципом нормировать всю свободную энергию: энтропию
состояния, мостовой детерминант, фермионный пфаффиан и классическое
действие. Делить только стабилизирующий член запрещено.

\section{Дискретный запрет конечных гауссовых мер}

Обычные конечномерные гауссовы интегралы дают степени детерминанта в
решётке
\begin{equation}
 \nu\in\frac12\mathbb Z:
\end{equation}
вещественный бозон даёт половинную степень, комплексный бозон --- целую,
а вещественный или комплексный фермион меняет знак соответствующей
степени. После любых конечных БВ/БРСТ-сокращений показатель остаётся в той
же решётке. Но
\begin{equation}
 \frac{17}{168}<\frac12.
\end{equation}
Следовательно, никакой ненулевой обычный конечный гауссов детерминант с
целой кратностью и линейным семейным блоком не попадает в найденное окно.

\begin{nogocriterion}[Запрет скрытого деления на размерность]
Нельзя выводить $1/20$ или $1/14$, сначала вычисляя обычный пфаффиан, а
затем деля только его логарифм на размерность удобного носителя. Такой шаг
равносилен новой аксиоме интенсивной свободной энергии. Она должна быть
применена ко всем секторам до повторного вычисления условия
неустойчивости.
\end{nogocriterion}

\section{Вердикт}

\begin{protocol}
\begin{enumerate}
 \item сумма KO-степеней внешнего и внутреннего факторов:
 \textbf{$4+6\equiv2\pmod8$};
 \item киральный product-KO2 пфаффиан:
 \textbf{кинематически допустим};
 \item явный семейный антисимметричный блок:
 \textbf{построен};
 \item обычный пфаффиановый показатель:
 \textbf{$15$};
 \item влияние множителя $1/300$ в квадратичном действии:
 \textbf{только $R$-независимая константа};
 \item нормированный показатель Фугледе--Кадисона:
 \textbf{$1/20$};
 \item стандартный фермионный интеграл выводит $1/20$:
 \textbf{нет};
 \item различие $M_{210}$ и $M_{300}$ выбирает обычную меру:
 \textbf{нет};
 \item ненулевой конечный гауссов показатель внутри окна:
 \textbf{невозможен};
 \item общий принцип интенсивной свободной энергии:
 \textbf{ещё не выведен};
 \item рождение материи:
 \textbf{не закрыто}.
\end{enumerate}
\end{protocol}

Следующая проверка должна проверить, допускает ли весь родитель единый
интенсивный функционал свободной энергии, в котором один и тот же
нормированный след применяется к классическому, мостовому, фермионному и
энтропийному секторам. Если прежняя отрицательная мода не переживает такую
общую перенормировку, детерминантная ветвь закрывается.

Машинный сертификат сохранён в
\path{s2t_v6_product_ko2_family_pfaffian_operator_gate_results.json}.